\section{Discussion and Limitations}
\label{sec:discussion}

The factor-separated evaluation supports the proposed mechanism rather than a
simple increase in planning effort. Adaptive computation uses fewer samples,
adaptive objectives reduce force error, and their combination raises compound
pass by 11.11 percentage points with an interval that excludes zero. The
improvement is concentrated in tracking and is especially visible at 12~N,
where strict tracking rises from 5/45 to 17/45. It does not increase task
completion or reduce sampled-limit violations. High-force regulation on the
cylinder therefore remains the central performance limit.

The mechanism ablations also delimit the contribution. Moving the historical
routing threshold changes how often MPPI is used but does not expose a natural
switching point. Conversely, adding a transient-envelope head or its
target-specific calibration does not improve the frozen development outcomes.
The supported method is consequently the smaller combination of persistent
force conditioning, learned next-force risk, continuous regime-conditioned
costs, and adaptive compute allocation.

The BC ablation clarifies what \emph{direct} control requires in this study.
The evaluation controller contains no teacher or separate force regulator,
but the actor proposal is learned with behavioural supervision. Removing that
supervision leaves one-step force-model accuracy unchanged and nevertheless
eliminates contact in every development run. The result attributes BC's main
role to proposal quality for contact acquisition and recovery, not to better
force prediction. The moderate coefficient further shows that low offline
action error is insufficient: its seed-dependent failures appear only in the
closed loop.

Planning latency remains a deployment constraint. Dynamic compute reduces the
mean sample count and host planning time in the development screen, but the
simulation still runs at a 100-Hz control discretization rather than a
demonstrated 100-Hz wall-clock planner. The method should therefore be read as
a control and allocation mechanism, not as evidence of real-time execution on
the present host.

The deployment-gap panel separates several practical sensitivities.  A
0.10-mm TCP-position error and 0.035-N RMS force noise leave aggregate
compound pass close to nominal, whereas a positive force bias sharply worsens
tracking in B0 and S1.  Fifty milliseconds of observation delay reduces both
completion and sampled-limit safety, and a 5-degree registered-normal error
primarily disrupts task completion.  The opposite-signed force bias sometimes
improves the thresholded tracking count while increasing force-tail exposure,
so calibration error cannot be summarized by a single success rate.  These
isolated interventions locate deployment sensitivities without representing
a joint hardware-error distribution.

The cross-engine results separate portability from task equivalence. The
revised MuJoCo port meets the held-out multilevel normal-cycle criteria and
matches the bidirectional motion's force history, onset, and peak, but its
0.833-mm tangential-position RMSE exceeds the 0.5-mm tolerance. Zero-shot
deployment is consequently interpreted as a stress test. All 30 flat and
inclined evaluations pass, whereas all 15 cylindrical evaluations fail with
sampled-limit exceedances. The surface-specific split strengthens the
evidence for the simpler geometries, but it does not establish engine-invariant
performance. Curved contact remains sensitive to the policy, geometry
representation, and solver dynamics together.

The repeated cylinder boundary also identifies a distinct method problem.
Prior curved-surface wiping and polishing systems couple local-normal
estimation, tool-posture adaptation, and force feedback
\cite{armesto2018constraintaware,li2023unknownsurface}.
Our policy receives a registered local frame but commands translation only;
neither its latent force predictor nor its adaptive MPPI objective can rotate
the tool to follow a changing normal. A curvature-specific extension should
therefore add a bounded posture channel and compare registered, force-derived,
and learned local-normal representations under a radius sweep. Retuning the
present scalar force costs would not isolate this missing capability. We leave
that extension outside the frozen study rather than adapting the method to the
observed cylinder blocks.

The numerical audit applies unchanged. Task completion, force tracking,
contact continuity, and sampled force tails have different evidence
requirements. The final method comparison consequently reports them
separately as well as jointly. Peak force is an empirical statistic under the
recorded solver configuration. The adaptive planner's mean peak increases by
0.247~N and both planners record two samples above 15~N, so the measured
tracking benefit is not interpreted as a safety improvement. MRE and NRMSE
describe regulation on the observed contact samples.

All experiments use rigid-contact simulation and a registered nominal task
frame. The observations are causal but do not reproduce a physical camera,
force-sensor bandwidth, calibration drift, deformable contamination, or robot
communication stack. The factor-separated design improves causal
interpretation within simulation. The cross-engine stress test broadens that
evidence without replacing hardware validation.
